\section{2. Related Work}

\subsection{2.1 Traditional Light Field Depth Estimation Methods}

\subsubsection{2.1 Traditional Light Field Depth Estimation Methods}

Many depth estimation methods for light-field cameras have been proposed, primarily leveraging Epipolar Plane Images (EPIs), defocus cues, and Multi-View Stereo (MVS) techniques to extract geometric information. Among these, EPI-based approaches explicitly extract geometric cues by analyzing the slope of linear structures formed by corresponding pixels across different views. Wanner et al. [?] proposed a robust depth estimation framework that derives depth from the orientation of EPI lines using structure tensors. Specifically, they formulated the depth estimation as an optimization problem with global smoothness constraints, successfully achieving state-of-the-art performance on Lambertian scenes. However, this method fundamentally relies on the strict photo-consistency assumption across viewpoints. When encountering non-Lambertian regions, such as specular reflections, the linear structures in EPIs are severely distorted or broken, leading to erroneous slope calculations and dramatic depth estimation failures.

While EPI-based methods primarily focus on angular correspondences, other approaches extend the geometric cue extraction by incorporating defocus cues to address the under-constrained nature of light-field depth estimation. Tao et al. [?] proposed a depth estimation method that synergistically combines defocus and correspondence cues. By analyzing the focal stack derived from the light field, they demonstrated that defocus cues provide complementary depth information, especially in textureless regions where correspondence cues fail. Although this fusion strategy remarkably improves depth accuracy in challenging Lambertian areas, it still struggles with spatially-varying reflectance properties. Specifically, the defocus blur of glossy or transparent materials behaves differently from that of diffuse surfaces, violating the implicit thin-lens defocus model and causing significant depth artifacts in mixed real-world scenes.

In contrast to EPI and defocus-based techniques, Multi-View Stereo (MVS) methods treat the light field as a dense array of sub-aperture images and leverage conventional stereo matching algorithms. To mitigate the impact of non-Lambertian reflections, Kim et al. [?] proposed an improved MVS framework tailored for light fields. They introduced a view-selection mechanism and a robust cost aggregation strategy to explicitly filter out specular highlights and outlier views, thereby preserving the photo-consistency of the remaining diffuse components. While this approach successfully generalizes MVS to certain glossy surfaces, it inherently treats non-Lambertian reflections as noise to be discarded rather than physical signals to be modeled. Consequently, it cannot handle general glossy surfaces or transparent materials where the diffuse component is absent or severely attenuated, inevitably resulting in missing depth values in highly reflective regions.

In short, although the aforementioned traditional methods have significantly advanced light-field depth estimation, they share critical limitations rooted in their physical assumptions. On the one hand, they heavily rely on the single Lambertian assumption. When encountering non-Lambertian regions such as specular or scattering materials, the viewpoint consistency is inevitably violated, leading to complete depth estimation failure or the notorious "texture copying" artifacts. On the other hand, these methods lack a unified physical modeling of light-matter interactions. They cannot explain the root cause of non-Lambertian failure from the perspective of physical optics (e.g., the rank variation of the Radiance Tensor Field), making it exceedingly difficult to effectively process urban mixed scenes containing diverse reflection types. In contrast, our proposed unified dual-mask physical model explicitly formulates the light transport process to accommodate pure diffuse, specular/scattering, and complex mixed scenes. Furthermore, by incorporating an MRI-inspired pixel-level angular frequency material parsing mechanism and a component-aware adaptive depth fusion strategy, our approach fundamentally addresses the physical root causes of non-Lambertian degradation, enabling robust and accurate depth recovery without relying on complex neural networks.

\subsection{2.2 Deep Learning-based General Light Field Depth Estimation Methods}

\subsubsection{2.2 Deep Learning-based General Light Field Depth Estimation Methods}

While traditional methods explicitly extract geometric cues based on strict photo-consistency assumptions, deep learning-based approaches have emerged to implicitly learn spatial-angular representations from light field data, demonstrating remarkable robustness in under-constrained scenarios. These general data-driven methods primarily leverage Convolutional Neural Networks (CNNs), 3D/4D convolutions, or Transformer architectures to regress depth maps end-to-end.

One of the pioneering works in this domain is EPINet, proposed by Shin et al. [?], which formulates depth estimation as a multi-directional Epipolar Plane Image (EPI) feature extraction problem. Specifically, EPINet extracts EPIs along four distinct angular directions and employs a 2D CNN to learn spatial-angular features, which are subsequently fused for depth regression. This approach successfully achieves state-of-the-art performance on standard Lambertian datasets by implicitly capturing the slope information of EPI lines. However, since EPINet and its variants (e.g., EPINet4Dir) fundamentally rely on the underlying linear structures of EPIs, they inevitably encounter severe performance degradation when the photo-consistency assumption is violated by specular reflections or transparent materials.

To capture the intrinsic high-dimensional structure of light fields more comprehensively, Heber et al. [?] introduced LFNet, which extends the feature extraction from 2D EPIs to the full 4D light field tensor. By utilizing 3D and 4D convolutions, LFNet explicitly models the joint spatial-angular correlations, significantly enlarging the receptive field and improving depth accuracy in occluded regions. Moreover, subsequent works have extensively refined this paradigm by incorporating multi-scale feature fusion and cost volume regularization. Despite these architectural advancements, processing the dense 4D tensor incurs dramatically high computational costs. More importantly, these 3D/4D CNNs remain highly sensitive to angular intensity variations; when encountering spatially-varying non-Lambertian materials, the network struggles to distinguish between geometric disparities and reflectance-induced intensity changes.

More recently, to address the limited receptive field inherent in CNNs, Transformer-based architectures have been introduced to light field depth estimation. For instance, the Light Field Transformer (LFT) [?] leverages self-attention mechanisms to model long-range spatial and angular dependencies globally. By treating light field patches as sequences, LFT successfully captures global context and achieves superior generalization on complex real-world scenes compared to purely convolutional models. Unfortunately, the attention mechanism tends to assign high weights to salient texture regions regardless of their material properties. Consequently, in non-Lambertian areas where angular consistency is broken, the Transformer implicitly correlates mismatched textures across views, leading to structurally plausible but physically incorrect depth predictions.

Despite their remarkable success on Lambertian datasets, these purely data-driven black-box models share critical limitations when encountering non-Lambertian regions. On the one hand, lacking explicit physical optical priors, these models are prone to overfitting to the texture distributions of the training set. When facing non-Lambertian regions, the network fails to decouple geometric disparity from reflectance variations, causing the prediction to degenerate into a mere "texture copying" of the input central view. On the other hand, complex deep models are highly susceptible to overfitting on the limited scale of available light field datasets. Furthermore, their standard loss functions optimize the global error uniformly, failing to adaptively weight the optimization based on pixel-level material confidence. As a result, the Mean Absolute Error (MAE) in non-Lambertian regions remains persistently high. In contrast to these general data-driven methods, our work introduces a unified dual-mask physical model and an MRI-inspired frequency-domain material parsing mechanism, explicitly addressing the physical root causes of depth failure in non-Lambertian scenes without relying on heavy black-box networks.

\subsection{2.3 Deep Learning-based Material-Aware and Non-Lambertian Processing Methods}

\subsubsection{2.3 Deep Learning-based Material-Aware and Non-Lambertian Processing Methods}

To address the severe performance degradation of general methods in non-Lambertian scenarios, recent studies have explicitly introduced material awareness, Bidirectional Reflectance Distribution Function (BRDF) estimation, or dual-cue mechanisms into deep light field processing pipelines. These approaches attempt to mitigate the violation of photo-consistency by distinguishing reflectance types or leveraging alternative geometric cues.

One representative direction is the integration of defocus cues with material features to compensate for the failure of Epipolar Plane Image (EPI)-based stereo matching. For instance, Sheng et al. [?] proposed MaterialDualCueNet, a dual-branch architecture that simultaneously extracts defocus and material-specific features. Specifically, the network employs a material-aware attention module to dynamically re-weight the defocus and EPI features based on the predicted surface reflectance. While this dual-cue strategy remarkably improves depth accuracy on moderately glossy surfaces, it fundamentally relies on the validity of defocus cues. Unfortunately, in highly specular or transparent regions, the defocus blur itself becomes spatially-varying and physically inconsistent, causing the defocus branch to extract misleading features and ultimately leading to severe depth artifacts.

While dual-cue networks focus on feature-level fusion, another line of research attempts to explicitly classify material types as an auxiliary step to guide depth regression. Chen et al. [?] introduced a data-driven, lightweight 3-class CNN to categorize pixels into diffuse, specular, and transparent regions. By training this classifier on central sub-aperture images and EPI slices, the method applies distinct heuristic depth refinement rules for each material category. Although this explicit classification successfully prevents the "texture copying" effect in purely specular regions to some extent, the 3-class CNN struggles to capture fine-grained physical material characteristics. Moreover, under the sparse angular sampling (e.g., $9 \times 9$) of commercial light-field cameras, the limited angular resolution severely restricts the CNN's receptive field in the angular dimension, resulting in coarse and often erroneous material boundaries.

Going beyond simple classification, more advanced frameworks attempt to decouple illumination and geometry directly from the light field via inverse rendering. Wang et al. [?] formulated an inverse rendering deep network that jointly estimates depth, surface normals, and intrinsic reflectance components (diffuse and specular). By incorporating a differentiable rendering layer, the network is constrained to reconstruct the input light field, thereby implicitly learning to separate shading from geometry. This physically-inspired formulation demonstrates impressive generalization on synthetic datasets and has been extensively evaluated on benchmark scenes. However, the decoupling process is entirely driven by implicit network learning rather than explicit physical modeling. Consequently, the multi-component fusion strategy relies on simple heuristic concatenation or learned weights, failing to adaptively leverage the intrinsic physical differences between various reflection types.

In short, while existing material-aware and non-Lambertian processing methods have made significant strides, they share critical common limitations that hinder their robustness in complex real-world scenes. On the one hand, existing learning-based material classifiers (e.g., CNNs) struggle to capture fine-grained physical material features under sparse angular sampling (e.g., $9 \times 9$), and the auxiliary defocus cues inevitably fail in strictly non-Lambertian regions where no valid geometric features can be extracted. On the other hand, these methods lack a pixel-level physical wave-vector analysis mechanism; their multi-component fusion strategies are predominantly simple heuristic concatenations or implicit network learning. They fail to achieve adaptive confidence-weighted fusion from a physical essence perspective, such as exploiting the angular frequency spectrum differences among distinct reflection types. In contrast, our proposed approach fundamentally addresses these deficiencies. By establishing a unified dual-mask physical model, we explicitly explain the physical root causes of depth estimation failure in non-Lambertian regions. Furthermore, we introduce an MRI-inspired pixel-level angular frequency material parsing method that operates in $\mathcal{O}(HW)$ complexity without complex neural networks, and a component-aware adaptive fusion strategy that completely resolves the EPI texture-copying failure, significantly enhancing depth accuracy and robustness in both mixed and purely non-Lambertian scenarios.

While existing material-aware and dual-cue methods attempt to mitigate photo-consistency violations, they predominantly rely on complex neural architectures and heuristic priors, lacking a unified physical framework to explain non-Lambertian failures and struggling with fine-grained material parsing under sparse angular sampling. Inspired by these observations, we propose a unified dual-mask physical modeling approach that addresses the aforementioned limitations by integrating an MRI-inspired pixel-level angular frequency analysis for real-time material parsing, alongside a component-aware adaptive depth fusion strategy to fundamentally resolve texture-copying artifacts in non-Lambertian regions.